\documentclass[a4paper,10pt]{article}

% Package imports
\usepackage{latexsym}
\usepackage{xcolor}
\usepackage{float}
\usepackage{ragged2e}
\usepackage[empty]{fullpage}
\usepackage{wrapfig}
\usepackage{lipsum}
\usepackage{tabularx}
\usepackage{titlesec}
\usepackage{geometry}
\usepackage{marvosym}
\usepackage{verbatim}
\usepackage{enumitem}
\usepackage{fancyhdr}
\usepackage{multicol}
\usepackage{graphicx}
\usepackage{cfr-lm}
\usepackage[T1]{fontenc}
\usepackage[utf8]{inputenc}
\usepackage[french]{babel}
\usepackage{fontawesome5}

% Color definitions
\definecolor{darkblue}{RGB}{0,0,139}

% Page layout
\setlength{\multicolsep}{0pt}
\pagestyle{fancy}
\fancyhf{} % clear all header and footer fields
\fancyfoot{}
\renewcommand{\headrulewidth}{0pt}
\renewcommand{\footrulewidth}{0pt}
\geometry{left=0.8cm, top=0.5cm, right=0.8cm, bottom=0.8cm}
\setlength{\footskip}{5pt}

% Hyperlink setup
\usepackage[hidelinks]{hyperref}
\hypersetup{
    colorlinks=true,
    linkcolor=darkblue,
    filecolor=darkblue,
    urlcolor=darkblue,
}

% Custom box settings
\usepackage[most]{tcolorbox}
\tcbset{
    frame code={},
    center title,
    left=0pt,
    right=0pt,
    top=0pt,
    bottom=0pt,
    colback=gray!20,
    colframe=white,
    width=\dimexpr\textwidth\relax,
    enlarge left by=-2mm,
    boxsep=4pt,
    arc=0pt,outer arc=0pt,
}

% URL style
\urlstyle{same}

% Text alignment
\raggedright
\setlength{\tabcolsep}{0in}

% Section formatting
\titleformat{\section}{
  \vspace{-4pt}\scshape\raggedright\large
}{}{0em}{}[\color{black}\titlerule \vspace{-7pt}]

% Custom commands
\newcommand{\resumeItem}[2]{
  \item{
    \textbf{#1}{\hspace{0.5mm}#2 \vspace{-0.5mm}}
  }
}

\newcommand{\resumePOR}[3]{
\vspace{0.5mm}\item
    \begin{tabular*}{0.97\textwidth}[t]{l@{\extracolsep{\fill}}r}
        \textbf{#1}\hspace{0.3mm}#2 & \textit{\small{#3}}
    \end{tabular*}
    \vspace{-2mm}
}

\newcommand{\resumeSubheading}[4]{
\vspace{0.5mm}\item
    \begin{tabular*}{0.98\textwidth}[t]{l@{\extracolsep{\fill}}r}
        \textbf{#1} & \textit{\footnotesize{#4}} \\
        \textit{\footnotesize{#3}} &  \footnotesize{#2}\\
    \end{tabular*}
    \vspace{-2.4mm}
}

\newcommand{\resumeProject}[4]{
\vspace{0.5mm}\item
    \begin{tabular*}{0.98\textwidth}[t]{l@{\extracolsep{\fill}}r}
        \textbf{#1} & \textit{\footnotesize{#3}} \\
        \footnotesize{\textit{#2}} & \footnotesize{#4}
    \end{tabular*}
    \vspace{-2.4mm}
}

\newcommand{\resumeSubItem}[2]{\resumeItem{#1}{#2}\vspace{-4pt}}

\renewcommand{\labelitemi}{$\vcenter{\hbox{\tiny$\bullet$}}$}
\renewcommand{\labelitemii}{$\vcenter{\hbox{\tiny$\circ$}}$}

\newcommand{\resumeSubHeadingListStart}{\begin{itemize}[leftmargin=*,labelsep=1mm]}
\newcommand{\resumeHeadingSkillStart}{\begin{itemize}[leftmargin=*,itemsep=1.7mm, rightmargin=2ex]}
\newcommand{\resumeItemListStart}{\begin{itemize}[leftmargin=*,labelsep=1mm,itemsep=0.5mm]}

\newcommand{\resumeSubHeadingListEnd}{\end{itemize}\vspace{2mm}}
\newcommand{\resumeHeadingSkillEnd}{\end{itemize}\vspace{-2mm}}
\newcommand{\resumeItemListEnd}{\end{itemize}\vspace{-2mm}}
\newcommand{\cvsection}[1]{%
\vspace{2mm}
\begin{tcolorbox}
    \textbf{\large #1}
\end{tcolorbox}
    \vspace{-4mm}
}

\newcolumntype{L}{>{\raggedright\arraybackslash}X}%
\newcolumntype{R}{>{\raggedleft\arraybackslash}X}%
\newcolumntype{C}{>{\centering\arraybackslash}X}%

% Commands for icon sizing
\newcommand{\socialicon}[1]{\raisebox{-0.05em}{\resizebox{!}{1em}{#1}}}
\newcommand{\ieeeicon}[1]{\raisebox{-0.3em}{\resizebox{!}{1.3em}{#1}}}

% Font options
\newcommand{\headerfontiii}{\fontfamily{ppl}\selectfont} % Palatino

\begin{document}
\headerfontiii

% Header
\begin{center}
    {\Huge\textbf{DKHILI FARES}}\\[2mm]
    {\large\textit{Ing\'enieur Syst\`emes Embarqu\'es | STM32 \& IoT}}
\end{center}
\vspace{-4mm}

\begin{center}
    \small{
    +216 27 491 238\,|\ ElManar-Tunis, Tunisie
    }
\end{center}
\vspace{-6mm}
\vspace{2.2mm}
\begin{center}
    \small{
    \href{mailto:Dkhili.fares2002@gmail.com}{Dkhili.fares2002@gmail.com} \,|\,
    \socialicon{\faLinkedin}
    \href{https://www.linkedin.com/in/fares-dkhili-018b17303/}{Fares Dkhili} \,|\,
    \socialicon{\faGithub}
    \href{https://github.com/DkhiliFares}{Fares Dkhili}
    }
\end{center}
\vspace{-2mm}
\begin{center}
\small{\textit{Ing\'enieur en syst\`emes embarqu\'es avec une exp\'erience pratique sur STM32 et MCU multi-vendeurs. Benchmarking des stacks middleware STM32Cube face \`a 5 \'ecosyst\`emes concurrents chez STMicroelectronics. Conception autonome d'outils d'automatisation pilot\'es par l'IA pour acc\'el\'erer l'analyse firmware.}}
\end{center}


\section{\textbf{Formation}}
\noindent
\textbf{Cycle Ing\'enieur en Informatique,} Facult\'e des Sciences de Tunis \hfill 2023 -- Pr\'esent\\
\textit{\footnotesize{Cours cl\'es : Syst\`emes Temps R\'eel, Architecture des Ordinateurs, Syst\`emes Embarqu\'es, Conception Num\'erique, R\'eseaux}}\\
\textbf{Cycle Pr\'eparatoire Int\'egr\'e en Math\'ematiques et Informatique,} Facult\'e des Sciences de Tunis \hfill 2021 -- 2023

\section{\textbf{Exp\'erience Professionnelle}}
\resumeSubHeadingListStart
\resumeSubheading
  {STMicroelectronics}{Tunis, Tunisie}
  {Stage de Fin d'\'Etudes (PFE) --- Benchmark \'Ecosyst\`eme STM32Cube}{F\'ev 2026 -- Pr\'esent}
  \resumeItemListStart
    \item Benchmark des stacks middleware STM32Cube (\textbf{USB, FreeRTOS, FatFS}) face \`a \textbf{5 \'ecosyst\`emes concurrents} (TI, NXP, Renesas, GigaDevice) en mesurant performance et empreinte m\'emoire
    \item Conception autonome d'un \textbf{analyseur d'empreinte pilot\'e par l'IA} avec interface chatbot --- automatise l'analyse des \'ecarts, le scoring et les recommandations d'optimisation sur les builds firmware MCU
    \item D\'eveloppement d'un \textbf{agent Copilot VS Code personnalis\'e} automatisant la cr\'eation de projets, le portage multi-vendeurs, la configuration de build et le d\'ebogage
    \item Pr\'esentation des r\'esultats d'analyse concurrentielle \`a l'\textbf{\'equipe marketing de STMicroelectronics}, identifiant les forces et axes d'am\'elioration par vendeur
    \item Outils d\'evelopp\'es en \textbf{totale autonomie} sans demande du superviseur, d\'emontrant initiative et maturit\'e technique
  \resumeItemListEnd

\resumeSubheading
  {STMicroelectronics}{Tunis, Tunisie}
  {Stage d'\'Et\'e --- Portage EEMBC CoreMark-PRO sur STM32}{Juillet -- Ao\^ut 2025}
  \resumeItemListStart
    \item R\'esolution d'une \textbf{d\'ependance POSIX bloquante} en portant le benchmark CoreMark-PRO de Linux vers STM32 bare-metal, adaptation des couches d'abstraction MITH pour ARM Cortex-M
    \item Ex\'ecution et comparaison des performances sur \textbf{Cortex-M4 et Cortex-M7} sur 3 cartes STM32 (H747I-EVAL, Nucleo-144, Discovery U5)
    \item Optimisation des charges m\'emoire limit\'ees via \textbf{SDRAM externe} et ajustement des configurations de build IAR/STM32CubeMX
  \resumeItemListEnd

\resumeSubheading
  {Vicom Consulting}{Tunis, Tunisie}
  {Stage d'\'Et\'e --- D\'eveloppement Linux Embarqu\'e}{Ao\^ut 2024}
  \resumeItemListStart
    \item Cross-compilation d'une \textbf{application Qt avec interface graphique} pour Raspberry Pi via une cha\^ine de compilation ARM et d\'eploiement sous Linux embarqu\'e
    \item Configuration du syst\`eme de build pour la \textbf{compatibilit\'e multiplateforme} et validation de l'ex\'ecution sur cible via SSH
  \resumeItemListEnd

\resumeSubHeadingListEnd
\vspace{-6mm}

\section{\textbf{Projets}}
\vspace{-0.4mm}
\resumeSubHeadingListStart

\resumeProject
  {Analyseur d'Empreinte Firmware pilot\'e par l'IA}
  {Outils : Python, Int\'egration LLM/IA, Interface Chatbot, G\'en\'eration de Rapports}
  {2026}
  {}
\resumeItemListStart
  \item Conception d'un outil \textbf{automatisant l'analyse d'empreinte firmware MCU} --- classifie les sections binaires, identifie les causes d'\'ecarts de taille entre vendeurs et g\'en\`ere des recommandations d'optimisation
  \item Int\'egration d'une \textbf{interface chatbot IA} pour l'interrogation interactive des donn\'ees d'empreinte et des r\'esultats de comparaison multi-vendeurs
\resumeItemListEnd

\resumeProject
  {Agent Copilot Benchmark pour STM32}
  {Outils : VS Code Agent API, TypeScript, STM32CubeMX, Cha\^ines multi-vendeurs}
  {2026}
  {}
\resumeItemListStart
  \item Construction d'un \textbf{agent Copilot VS Code personnalis\'e} automatisant l'ensemble du workflow benchmark : cr\'eation, portage multi-vendeurs, configuration de build et setup de d\'ebogage
\resumeItemListEnd

\resumeProject
  {SecureAccess --- Syst\`eme d'Authentification Multi-Facteurs}
  {Outils : Python, Raspberry Pi, OpenCV, Flask, FaceRecognition, RFID}
  {Mai 2025}
  {{}\href{https://github.com/DkhiliFares/SecureAccess}{\textcolor{darkblue}{\faGithub}}}
\resumeItemListStart
  \item Conception d'un syst\`eme de s\'ecurit\'e physique combinant \textbf{3 facteurs d'authentification} (RFID, reconnaissance faciale, code PIN) avec tableau de bord Flask
\resumeItemListEnd

\resumeProject
{Cross-Compilation Linux Embarqu\'e}
{Outils : Buildroot, GCC, U-Boot, Raspberry Pi, BeagleBone Black}
{Novembre 2025}
{{}\href{https://github.com/DkhiliFares/Embedded-linux-cross-compilation}{\textcolor{darkblue}{\faGithub}}}
\resumeItemListStart
\item Configuration d'une \textbf{cha\^ine de cross-compilation} et personnalisation d'une image Linux embarqu\'e minimale avec Buildroot pour plateformes ARM
\resumeItemListEnd

\resumeProject
  {Biblioth\`eque de Communication GPRS pour IoT}
  {Outils : STM32, C, SIM808, UART}
  {D\'ecembre 2024}
  {{}\href{https://github.com/DkhiliFares/GPRS-DRIVER}{\textcolor{darkblue}{\faGithub}}}
\resumeItemListStart
  \item Impl\'ementation d'une \textbf{biblioth\`eque HAL STM32 r\'eutilisable} pour le module SIM808, abstrayant les commandes AT en APIs HTTP GET/POST propres via GPRS
\resumeItemListEnd

\resumeProject
  {Infrastructure d'Administration R\'eseau}
  {Outils : EVE-NG, Cisco IOL, Zabbix, Akvorado, Wireshark}
  {Mars 2025}
  {{}\href{https://github.com/DkhiliFares/Network_Administration_Project}{\textcolor{darkblue}{\faGithub}}}
\resumeItemListStart
  \item D\'eploiement d'une \textbf{topologie multi-VLAN entreprise} sur EVE-NG avec supervision centralis\'ee (Zabbix) et analyse de trafic (Akvorado)
\resumeItemListEnd

\resumeProject
  {Syst\`eme IoT de D\'etection de Gaz et Incendie}
  {Outils : ESP8266, C++, DHT11, MQ2, MQTT}
  {Mai 2024}
  {{}\href{https://github.com/DkhiliFares/FireDetection__esp32}{\textcolor{darkblue}{\faGithub}}}
\resumeItemListStart
  \item R\'ealisation d'un \textbf{n\oe ud capteur IoT} pour la d\'etection en temps r\'eel de gaz/fum\'ee avec alertes imm\'ediates via tableau de bord cloud
\resumeItemListEnd
\resumeSubHeadingListEnd

\section{\textbf{Comp\'etences}}
\textbf{Programmation :} C, C++ (POO), Embedded C, Python, VHDL, Bash, Java, TypeScript \\
\textbf{Embarqu\'e \& RTOS :} STM32 (CubeMX, CubeIDE, Keil, IAR), ESP32/8266, Raspberry Pi, FreeRTOS, Bare-metal, ARM Cortex-M, Pilotes p\'eriph\'eriques (GPIO, UART, SPI, I2C, CAN) \\
\textbf{Linux \& Syst\`emes de Build :} Buildroot, Yocto, Cross-compilation, Makefile/CMake, Scripts Shell, U-Boot \\
\textbf{IoT \& Connectivit\'e :} MQTT, LoRa/WAN, GPRS, Node-RED, Tableaux de bord Cloud, Power BI \\
\textbf{IA \& Automatisation :} Int\'egration LLM, Outillage pilot\'e par l'IA, OpenCV, D\'eveloppement d'agents VS Code \\
\textbf{FPGA \& Conception Num\'erique :} VHDL, Quartus Prime, ModelSim, Conception FSM, Co-design HW/SW \\
\textbf{DevOps \& Outils :} Git/GitHub, Docker, Jenkins, SonarQube, Pipelines CI/CD \\
\textbf{R\'eseaux \& S\'ecurit\'e :} VLAN, Routage, ACL, Cisco CLI (CCNA), Wireshark, Zabbix, TLS/SSL, VPN \\
\textbf{Langues :} Fran\c{c}ais (B2), Anglais (B1), Arabe (Natif)

\section{\textbf{Certifications}}
\vspace{-0.2mm}
\resumeSubHeadingListStart
\resumePOR{}{\href{https://www.linkedin.com/in/fares-dkhili-018b17303/recent-activity/all/}{
\textbf{CCNA : Introduction aux R\'eseaux (Cisco)}
}}{Octobre 2024}
\resumePOR{}{\href{https://learn.nvidia.com/certificates?id=Z4jS_G2wT_uK63hkrrwHAw}{
\textbf{Fundamentals of Deep Learning (NVIDIA)}
}}{Avril 2025}
\resumeSubHeadingListEnd
\vspace{-6mm}

\section{\textbf{Distinctions}}
\vspace{-0.4mm}
\resumeSubHeadingListStart
\resumeProject
  {Performance en Comp\'etition}
  {CTF 101, Securinet FST}
  {Avril 2025}
  {{}\href{https://www.linkedin.com/posts/fares-dkhili-018b17303_cybersecurity-ctf-activity-7322699787982528512-YMcf?utm_source=share&utm_medium=member_desktop&rcm=ACoAAE2V9FQBl97M7EPoPtI-0a16X0hT_ahlsmk}{\textcolor{darkblue}{\faIcon{globe}}}}
\resumeItemListStart
  \item \textbf{3\`eme place} au CTF 101 (Securinet FST) --- r\'esolution de challenges de cryptographie et forensique sous contrainte de temps
\resumeItemListEnd
\resumeSubHeadingListEnd

\vspace{-6mm}
\section{\textbf{Engagements Associatifs}}
\vspace{-0.4mm}
\resumeSubHeadingListStart
\resumePOR{Securinets FST Tunis}
    {}
    {Novembre 2024 -- Mai 2025}
\resumePOR{ATIA FST : Club IA}
    {}
    {2025 -- Pr\'esent}
\resumePOR{GDG on Campus : FST}
    {}
    {Septembre 2023 -- Pr\'esent}
\resumeSubHeadingListEnd
\vspace{-6mm}
\end{document}
